\documentclass[11pt]{article}

% ---- ICLR-preprint-like, self-contained preamble ----
\usepackage[letterpaper, margin=1in]{geometry}
\usepackage{amsmath,amssymb,amsthm,mathtools}
\usepackage{graphicx}
\usepackage{booktabs}
\usepackage{microtype}
\usepackage{xcolor}
\usepackage[colorlinks=true, linkcolor=blue!60!black, citecolor=blue!60!black, urlcolor=blue!60!black]{hyperref}
\usepackage[numbers,sort&compress]{natbib}
\usepackage{caption}
\captionsetup{font=small, labelfont=bf}
\usepackage{titlesec}
\titleformat*{\section}{\large\bfseries}
\titleformat*{\subsection}{\normalsize\bfseries}
\setlength{\bibsep}{2pt}

\newcommand{\eg}{\varepsilon_g}
\newcommand{\R}{\mathbb{R}}
\newcommand{\E}{\mathbb{E}}
\newcommand{\KL}{\mathrm{KL}}
\newcommand{\bx}{\boldsymbol{x}}
\newcommand{\bmu}{\boldsymbol{\mu}}
\newcommand{\btheta}{\boldsymbol{\theta}}

\title{\bf Phase Diagrams Without Solvable Models:\\
An Empirical Teacher--Student Atlas of Modern Learning Paradigms}

\author{%
Yuma Ichikawa\\
Fujitsu Limited / RIKEN\\
\texttt{yuma.ichikawa@a.riken.jp}
}
\date{}

\begin{document}
\maketitle

\begin{abstract}
The statistical mechanics of learning explains generalization in neural networks through exactly solvable teacher--student models, whose order parameters reveal phase transitions invisible to worst-case analysis. Modern learning paradigms---supervised fine-tuning, learning from reward models, weak-to-strong supervision, training on synthetic data, and in-context learning---lie far beyond the reach of current theory, yet there is no reason to believe they are free of such collective phenomena. In this work we take the experimental route: we formulate each paradigm as a minimal nonlinear teacher--student problem, track the same macroscopic order parameters that theory uses (teacher--student overlap, generalization error, and paradigm-specific observables), and map empirical phase diagrams over the natural control parameters. Across all five paradigms we find sharp, physics-like structure: a similarity--data boundary separating transfer from catastrophic forgetting, a Goodhart point where reward optimization turns against the true objective, a supervisor-accuracy threshold for weak-to-strong gains, a real-data fraction that arrests model collapse, and a task-diversity transition from memorization to in-context generalization. A teacher taxonomy---random, structured, and genuinely trained on real images---shows that these boundaries shift systematically with teacher structure. All settings are released as a documented, extensible open-source framework, offering theory a set of concrete, quantitative targets.
\end{abstract}

\section{Introduction}
\label{sec:intro}

Understanding why large learning systems generalize remains a central open problem, and the statistical mechanics of learning has long offered one of the few quantitative windows into it. In the teacher--student framework, a network learns from data labeled by a hidden ``teacher'' network, and in the high-dimensional limit the entire learning problem collapses onto a handful of macroscopic \emph{order parameters}---the teacher--student overlap, the student self-overlap, and the generalization error---whose deterministic evolution can be computed exactly by the replica method or dynamical mean-field theory \citep{gardner1989three,seung1992statistical,engel2001statistical,zdeborova2016statistical}. This program has explained specialization transitions in shallow networks \citep{saad1995exact}, the double-descent behavior of overparametrized models \citep{belkin2019reconciling,nakkiran2021deep,mei2022generalization}, and, in a line of work steadily pushing minimal models toward realism, learning on structured data \citep{goldt2020modelling,gerace2020generalisation,loureiro2021learning} and the dynamics of stochastic training \citep{mignacco2020dynamical}.

The learning systems deployed today, however, are trained by pipelines that this theory has not yet reached. A modern model is pretrained, then fine-tuned on a narrower task (SFT); aligned by optimizing against a learned reward model (RLHF); supervised by weaker models (weak-to-strong generalization); increasingly trained on data generated by earlier models (a synthetic-data loop); and used through in-context learning rather than weight updates. Each of these paradigms composes multiple nonlinear learning stages, which places them beyond current exact analysis. Empirically, each is also known to exhibit abrupt, collective phenomena---catastrophic forgetting \citep{french1999catastrophic,kirkpatrick2017overcoming}, reward overoptimization \citep{gao2023scaling}, model collapse \citep{shumailov2024ai}, and the sudden emergence of in-context learning \citep{garg2022can,raventos2023pretraining}---that are strongly reminiscent of phase transitions in solvable models.

This mismatch defines the gap we address. The theory side owns exact order parameters but only for single-stage, effectively convex or narrowly structured problems; the empirical side owns the modern paradigms but studies them at full scale, where confounds (data curation, architecture, optimization tricks) obscure whatever collective structure is present, and where the quantities theory could one day predict are not even measured. What is missing is the intermediate object that statistical mechanics has always relied on: \emph{minimal models of the modern paradigms, instrumented with the order parameters of the theory}, small enough to sweep control parameters densely and cleanly, yet retaining the nonlinearity that makes the paradigms interesting.

In this work, we construct exactly this object. We formulate five modern paradigms---SFT, reward-model optimization, weak-to-strong supervision, recursive synthetic-data training, and in-context learning---as minimal teacher--student problems with nonlinear (tanh MLP or transformer) teachers and students. For each we identify the natural control parameters (task similarity, sample ratios $\alpha = n/d$, optimization strength, real-data fraction, task diversity) and paradigm-specific order parameters (function-space overlaps $\hat m$, proxy--gold reward gaps, performance gap recovered, output-variance ratios, in-context alignment), and we measure empirical phase diagrams by dense sweeps at moderate dimension ($d = 32$--$64$) with seed-averaged statistics. Orthogonally to the paradigms, we introduce a \emph{teacher taxonomy}---random weights, structured weights (sparse, low-rank, heavy-tailed, spiked, binary), and teachers genuinely trained on real image data---and cross every teacher with every paradigm, asking whether the observed boundaries are universal or structure-dependent (Figure~\ref{fig:taxonomy}).

Our experiments reveal sharp, reproducible structure in all five paradigms: (i) a transfer--forgetting boundary in the (task-similarity, data-ratio) plane of SFT; (ii) a Goodhart point $\KL^*$ at which optimizing a learned reward begins to \emph{decrease} the true reward, moving systematically with reward-model data; (iii) a supervisor-competence threshold above which a strong student reliably surpasses its weak teacher; (iv) an anchoring boundary in the real-data fraction that arrests model collapse; and (v) a task-diversity transition at which a transformer switches from memorizing pretraining tasks to implementing in-context regression. The taxonomy cross shows these boundaries are \emph{not} universal: e.g.\ the Goodhart point varies by a factor of $\sim 40$ between linear and low-rank teachers, and weak-to-strong gains that are robust for synthetic teachers vanish for teachers trained on real images.

Our contributions are: (1) a unified formulation of five modern learning paradigms as minimal nonlinear teacher--student problems with explicit, physics-style order parameters; (2) empirical phase diagrams for each paradigm, exposing sharp boundaries and their finite-size trends; (3) a teacher taxonomy and a teacher$\times$paradigm cross experiment quantifying which phenomena are structure-sensitive; and (4) an open-source, tested, documented framework in which every figure of this paper is one command, designed so new paradigms and teachers compose. We view these measurements as experimental targets for the theory that does not yet exist, in the same way that early simulations of spin glasses preceded their analytical understanding.

\section{Preliminaries: order parameters in teacher--student learning}
\label{sec:prelim}

\paragraph{Teacher--student setting.}
A teacher is a fixed function $f^*:\R^d \to \R$, here a nonlinear network with weights drawn from a prescribed ensemble or trained on real data. Data are generated as
\begin{equation}
    y_\mu = f^*(\bx_\mu) + \sigma \xi_\mu, \qquad \bx_\mu \sim P_x, \quad \xi_\mu \sim \mathcal{N}(0,1), \quad \mu = 1,\dots,n,
\end{equation}
and a student $f_{\btheta}$ is trained on $\{(\bx_\mu, y_\mu)\}$ by empirical risk minimization. The central macroscopic quantities of the solvable theory are the teacher--student overlap and the generalization error. Because our teachers and students are generic nonlinear networks, we measure both in \emph{function space}: with a fresh probe set $\{\bx_i\}_{i=1}^{N_p} \sim P_x$,
\begin{equation}
    \hat m \;=\; \frac{\langle f_{\btheta} f^* \rangle_p}{\sqrt{\langle f_{\btheta}^2\rangle_p\,\langle (f^*)^2\rangle_p}},
    \qquad
    \eg \;=\; \big\langle \big(f_{\btheta} - f^*\big)^2 \big\rangle_p ,
    \label{eq:overlap}
\end{equation}
where $\langle\cdot\rangle_p$ denotes the probe average and teacher outputs are noiseless. For linear models \eqref{eq:overlap} reduces to the classical weight-space overlap $m = \boldsymbol{w}\cdot\boldsymbol{w}^*/d$; for nonlinear models it remains well-defined, architecture-agnostic, and directly comparable across paradigms. The key control parameter throughout is the sample ratio $\alpha = n/d$.

\paragraph{Phenomenology of the solvable regime.}
In exactly solvable settings, $(\hat m, \eg)$ as functions of $\alpha$ display sharp thresholds (e.g.\ a critical $\alpha_c$ for weak-recovery in single-index models), specialization transitions with plateau dynamics \citep{saad1995exact}, and double descent at the interpolation threshold \citep{belkin2019reconciling,mei2022generalization}; structured data shift these boundaries in ways captured by hidden-manifold and Gaussian-universality analyses \citep{goldt2020modelling,gerace2020generalisation,loureiro2021learning}. Our working hypothesis is that the same observables, measured in the multi-stage paradigms below, will expose analogous structure---and that its \emph{location and sharpness} are the quantities a future theory must predict.

\section{Five paradigms as minimal teacher--student problems}
\label{sec:settings}

All settings share one toolbox (a tanh MLP $f(\bx) = \boldsymbol{a}^\top \tanh(W\bx/\sqrt{d})$ unless stated, one full-batch Adam training loop with early stopping, the probe-based observables \eqref{eq:overlap}); each paradigm adds only its protocol. Complete formulas and defaults are given in Appendix~\ref{app:settings}.

\subsection{SFT as a two-teacher problem}
\label{sec:set-sft}
Two teachers with weight-space correlation $\rho$ (task similarity) are built by interpolation, $\btheta_B^* = \rho\,\btheta_A^* + \sqrt{1-\rho^2}\,\btheta_\perp$. A student is pretrained on $n_{\rm pre} = \alpha_{\rm pre} d$ samples from teacher $A$, then fine-tuned on $n_{\rm ft} = \alpha_{\rm ft} d$ samples from teacher $B$, tracking $\hat m_A(t)$ and $\hat m_B(t)$. Order parameters: \emph{forgetting} $F = \hat m_A^{\rm pre} - \hat m_A^{\rm post}$ and \emph{transfer gain} $\Delta_B = \hat m_B^{\rm post} - \hat m_B^{\rm scratch}$, the latter against a from-scratch control with an identical budget. Control plane: $(\rho, \alpha_{\rm ft})$.

\subsection{Reward-model overoptimization (RLHF)}
\label{sec:set-rlhf}
A gold reward $r^*(\bx) = f^*(\bx) - \lambda\,\phi(\|\bx\|^2/d)$ combines a teacher network with an off-distribution penalty $\phi$ (hackability: high proxy regions off the data manifold carry low true reward). A proxy reward model $r_\phi$ is trained on $n_{\rm pref} = \alpha_r d$ Bradley--Terry preference pairs \citep{bradley1952rank} sampled in-distribution. A policy $\pi_{\bmu} = \mathcal{N}(\bmu, I)$ is then optimized by gradient ascent on $\E_{\pi}[r_\phi] - c\,\KL(\pi \| \pi_0)$ for a grid of KL penalties $c$. Order parameters: z-scored proxy and gold rewards along the optimization-strength axis $\KL(\pi\|\pi_0) = \|\bmu\|^2/2$, the \emph{Goodhart point} $\KL^*$ (gold-reward argmax), peak gold reward $G_{\max}$, and reward-model quality $m_{\rm RM}$. Control parameter: $\alpha_r$.

\subsection{Weak-to-strong generalization}
\label{sec:set-w2s}
A truth teacher generates $n_w = \alpha_w d$ labels for a \emph{weak} low-capacity supervisor; the weak model then labels $n_s = \alpha_s d$ fresh inputs for a \emph{strong} high-capacity student \citep{burns2023weak}. Order parameters: $\hat m_{\rm weak}, \hat m_{\rm strong}$ (both against the truth), imitation overlap $\hat m_{\rm imit}$, and the \emph{performance gap recovered} $\mathrm{PGR} = (\hat m_{\rm strong} - \hat m_{\rm weak})/(\hat m_{\rm ceiling} - \hat m_{\rm weak})$, with the ceiling trained on true labels. Control plane: $(\alpha_w, \alpha_s)$.

\subsection{Model collapse under recursive synthetic data}
\label{sec:set-collapse}
Generation $0$ trains on true teacher labels; generation $g{+}1$ trains on labels produced by generation $g$, with a fraction $p_{\rm real}$ of labels drawn from the true teacher \citep{shumailov2024ai}. Order parameters: overlap-with-truth $\hat m(g)$ and output-variance ratio $q(g)/q(0)$ (tail collapse). Control parameters: $(p_{\rm real}, \alpha)$.

\subsection{Emergence of in-context learning}
\label{sec:set-icl}
A small causal transformer is pretrained on regression sequences $(\bx_1, y_1, \dots, \bx_k, y_k, \bx_{k+1})$ whose tasks $\boldsymbol{w}$ are drawn from a finite pool of size $N_{\rm tasks}$ \citep{garg2022can,raventos2023pretraining}. Order parameters: ICL score on \emph{unseen} tasks, memorization score on pool tasks, and \emph{ridge alignment}---the correlation between the transformer's in-context predictions and the Bayes-optimal ridge predictor. Control parameter: $N_{\rm tasks}$.

\subsection{The teacher taxonomy}
\label{sec:set-tax}
Orthogonally to the paradigms, teachers are drawn from a registry spanning three families: \textbf{random} (i.i.d.\ Gaussian MLP; linear single-index), \textbf{structured} (sparse $90\%$-zero, low-rank, heavy-tailed, rank-one spiked, binary weights, all on Gaussian inputs; plus a random-weight MLP evaluated on real digit images), and \textbf{trained} (an MLP genuinely trained to $>90\%$ accuracy on digit-parity over real $8{\times}8$ images \citep{pedregosa2011scikit}, evaluated on the real image manifold). Each entry supplies $(f^*, P_x)$, so the data distribution and the teacher are injected jointly into any paradigm; the trained/random pair on identical real inputs isolates weight structure from data structure. The cross experiment probes every (teacher, paradigm) cell with the paradigm's headline order parameter.

\section{Experiments}
\label{sec:experiments}

\paragraph{Protocol.}
All experiments use $d = 32$--$64$, tanh MLP students of width $16$--$32$ (a $4$-layer transformer for ICL), full-batch Adam, and $N_p \in [2048, 4096]$ probe samples; sweeps average $2$--$5$ seeds and error bars are $\pm 1$ std. Every figure is reproducible with one CLI command (\texttt{statphys study \{sft,rlhf,weak\_to\_strong,collapse,icl,taxonomy\}}); raw JSON for every panel is emitted alongside the figure. We first present each paradigm's phase diagram, then the taxonomy cross.

\begin{figure}[t]
    \centering
    \includegraphics[width=\linewidth]{figures/sft.png}
    \caption{\textbf{SFT as a two-teacher problem} ($d{=}64$). \emph{Left:} fine-tuning trajectories $\hat m_A(t)$ (solid) and $\hat m_B(t)$ (dashed) at three task similarities $\rho$; forgetting of task $A$ is fast, monotone, and nearly total for $\rho \le 0.5$. \emph{Right:} forgetting $F$ in the $(\rho, \alpha_{\rm ft})$ plane. Forgetting depends dominantly on $\rho$, with a sharp drop above $\rho \approx 0.8$; the cyan contour marks the sign change of the transfer gain $\Delta_B$, separating a positive-transfer phase (high similarity, scarce fine-tuning data) from a regime where pretraining does not help.}
    \label{fig:sft}
\end{figure}

\subsection{SFT: a transfer--forgetting phase boundary}
Figure~\ref{fig:sft} shows the SFT plane. Forgetting is governed almost entirely by task similarity: $F \gtrsim 0.9$ for $\rho \lesssim 0.4$ at large $\alpha_{\rm ft}$, collapsing to $F \approx 0.05$ at $\rho = 1$. The dependence on $\alpha_{\rm ft}$ is comparatively weak---even $\alpha_{\rm ft} = 0.5$ erases a dissimilar task A ($F \approx 0.67$ at $\rho = 0$)---consistent with forgetting being an \emph{interference} effect rather than a capacity effect in this regime. Transfer, in contrast, has a genuine boundary: $\Delta_B > 0$ only above $\rho \approx 0.6$ and at small $\alpha_{\rm ft}$ (where the prior matters most), with $\Delta_B \approx +0.2$ at $(\rho, \alpha_{\rm ft}) = (1.0, 0.5)$. The trajectories (left) show the asymmetry of the two processes: learning B and forgetting A happen on the same fast timescale, after which both overlaps are static.

\begin{figure}[t]
    \centering
    \includegraphics[width=\linewidth]{figures/rlhf.png}
    \caption{\textbf{Reward-model overoptimization} ($d{=}32$). \emph{Left:} z-scored gold reward (solid) and proxy reward (dotted) along the optimization-strength axis $\KL(\pi\|\pi_0)$, for reward models trained with $\alpha_r = 1$--$64$ preference pairs per dimension. The proxy rises monotonically while the gold reward peaks and then falls---the Goodhart turnover. \emph{Right:} the Goodhart point $\KL^*$ and the peak gold value $G_{\max}$ both increase with $\alpha_r$ and saturate near $\alpha_r \approx 16$, while reward-model quality $m_{\rm RM}$ (blue diamonds) is non-monotone at large $\alpha_r$.}
    \label{fig:rlhf}
\end{figure}

\subsection{RLHF: the Goodhart point and its data dependence}
Figure~\ref{fig:rlhf} shows the overoptimization sweep. For every $\alpha_r$, the proxy reward increases monotonically with optimization strength while the gold reward reaches a maximum at a finite $\KL^*$ and then decreases steeply---at strong optimization the policy sits in regions where the proxy is confidently wrong, and the gold reward falls to $-8$ standard deviations. Better reward models postpone but do not remove the turnover: $\KL^*$ grows from $\approx 2.6$ ($\alpha_r = 1$) to $\approx 5$ ($\alpha_r \ge 16$) and then saturates, and the attainable utility $G_{\max}$ saturates similarly. This reproduces, in a $d = 32$ teacher--student model, the qualitative law observed at LLM scale by \citet{gao2023scaling}, and locates the regime (saturation of $\KL^*$ in $\alpha_r$) where more preference data stops buying more safely-optimizable reward.

\begin{figure}[t]
    \centering
    \includegraphics[width=\linewidth]{figures/weak_to_strong.png}
    \caption{\textbf{Weak-to-strong generalization} ($d{=}64$; weak width $2$, strong width $32$). \emph{Left:} $\hat m_{\rm strong}$ vs $\hat m_{\rm weak}$ across the $(\alpha_w, \alpha_s)$ sweep; points above the diagonal indicate students surpassing their supervisors, which occurs consistently for $\hat m_{\rm weak} \lesssim 0.9$. \emph{Right:} PGR in the $(\alpha_w, \alpha_s)$ plane: gains concentrate at intermediate supervisor competence ($\alpha_w \approx 4$--$10$), with noisy, even negative cells when the weak supervisor is already near ceiling.}
    \label{fig:w2s}
\end{figure}

\subsection{Weak-to-strong: gains require an imperfect but competent supervisor}
Figure~\ref{fig:w2s} shows the weak-to-strong surface. The strong student ends above the diagonal in essentially every cell---imitating a weak supervisor on ample data ($\alpha_s \gg \alpha_w$) yields $\hat m_{\rm strong} > \hat m_{\rm weak}$, i.e.\ the student averages out part of the supervisor's error rather than copying it. The recovered fraction is structured: PGR climbs from $\approx 0$ at $\alpha_w = 1$ (too weak a supervisor: nothing to build on) to $\approx 0.35$--$0.75$ at $\alpha_w = 4$--$10$, then becomes noisy and can turn negative as $\hat m_{\rm weak} \to \hat m_{\rm ceiling}$ (the denominator vanishes). The regime with reliable gains---an imperfect but competent supervisor and abundant weak labels---matches the operating point reported at scale \citep{burns2023weak}.

\begin{figure}[t]
    \centering
    \includegraphics[width=\linewidth]{figures/collapse.png}
    \caption{\textbf{Model collapse under recursive synthetic data} ($d{=}64$, $\alpha{=}8$). \emph{Left:} overlap with truth $\hat m(g)$ over $12$ generations for real-data fractions $p_{\rm real} \in [0,1]$: the fully synthetic loop ($p_{\rm real}=0$) decays steadily, while $p_{\rm real}\ge 0.3$ is statistically indistinguishable from fully real training. \emph{Middle:} output-variance ratio $q(g)/q(0)$; only the fully synthetic loop contracts. \emph{Right:} terminal overlap vs $p_{\rm real}$---an anchoring boundary at small $p_{\rm real}$.}
    \label{fig:collapse}
\end{figure}

\subsection{Model collapse: a small real-data anchor suffices}
Figure~\ref{fig:collapse} shows the recursive-training loop. With $p_{\rm real} = 0$ the overlap decays roughly linearly in generation number ($0.92 \to 0.86$ over $12$ generations at $\alpha = 8$) and the output variance contracts by $\sim 15\%$---both signatures of collapse \citep{shumailov2024ai}, here measured as clean order parameters. The striking feature is the smallness of the required anchor: terminal overlap recovers steeply in $p_{\rm real}$, with $p_{\rm real} \approx 0.3$ already indistinguishable from $p_{\rm real} = 1$ within error bars. In this regime collapse is not a cliff but a slow drift, and modest real-data replay arrests it---a quantitative target for theories of self-consuming loops.

\subsection{ICL: from task retrieval to in-context regression}
Figure~\ref{fig:icl} shows the task-diversity sweep. At small diversity the transformer is a task-retrieval system: memorization score $\approx 1$ on pool tasks, strongly negative ICL score on unseen tasks, and near-zero alignment with the ridge predictor ($N_{\rm tasks} \le 4$: ICL $= -0.94$ to $-0.34$, alignment $= 0.01$--$0.19$). The transition is fast and localized: between $N_{\rm tasks} = 8$ and $16$ the ICL score jumps from $0.10$ to $0.71$ and the ridge alignment from $0.47$ to $0.84$, after which both saturate ($\approx 0.86$ and $\approx 0.93$ at $N_{\rm tasks} = 1024$) while the memorization score decays only mildly---the network stops storing tasks and starts implementing a general in-context regression algorithm, reproducing the task-diversity threshold of \citet{raventos2023pretraining} as a two-order-parameter transition. The unseen-task error drops by an order of magnitude across the transition but remains well above the Bayes-optimal ridge error, quantifying how far this minimal transformer is from the optimal in-context algorithm it approximates.

\begin{figure}[t]
    \centering
    \includegraphics[width=\linewidth]{figures/icl.png}
    \caption{\textbf{Emergence of in-context learning} ($d{=}8$, context $k{=}16$, $2$ seeds). \emph{Left:} ICL score on unseen tasks (red), memorization score on pool tasks (blue), and alignment with the Bayes-optimal ridge predictor (orange) vs pretraining task diversity $N_{\rm tasks}$. The system switches from task retrieval to in-context regression between $N_{\rm tasks} = 8$ and $16$. \emph{Right:} unseen-task test error against the ridge error: the transition drops the error by an order of magnitude, without reaching the Bayes-optimal floor.}
    \label{fig:icl}
\end{figure}

\begin{figure}[t]
    \centering
    \includegraphics[width=\linewidth]{figures/taxonomy.png}
    \caption{\textbf{Teacher taxonomy $\times$ paradigm cross} ($2$ seeds; bars $\pm 1$ std). One headline order parameter per paradigm, for nine teachers spanning random (blue), structured (orange), and genuinely trained-on-real-images (red) families. Teacher structure shifts every boundary: the Goodhart point spans $\KL^* \approx 0.35$ (low-rank) to $\approx 17$ (linear); weak-to-strong gains vanish for the digits-trained teacher; collapse is fastest for sparse and trained teachers; forgetting is comparatively universal.}
    \label{fig:taxonomy}
\end{figure}

\subsection{The taxonomy cross: which phenomena are universal?}
Figure~\ref{fig:taxonomy} compiles the teacher$\times$paradigm table. Three findings stand out. \textbf{(i) Hackability tracks teacher complexity:} the linear teacher's preferences produce a proxy that stays faithful far off-distribution ($\KL^* \approx 16.6$), while the low-rank teacher is hacked almost immediately ($\KL^* \approx 0.35$)---a $\sim\!40\times$ shift of the safety-relevant boundary from weight structure alone. \textbf{(ii) Weak-to-strong gains are structure-fragile:} PGR is large for the linear teacher ($0.67$) and moderate for most structured ones ($\approx 0.5$), but consistently $\approx 0$ for the teacher trained on real images---its weak supervisor's errors are inherited, not averaged out; the comparison with its random-weight control on identical inputs ($\mathrm{PGR} = 0.14$) attributes the loss to \emph{trained weight structure}, not to the data manifold. \textbf{(iii) Collapse speed is teacher-dependent:} sparse and trained teachers lose signal fastest under the synthetic loop ($\Delta \hat m \approx 0.12$ in $8$ generations), while the spiked teacher barely decays ($\approx 0.01$): functions concentrated on few directions survive self-distillation. Forgetting, by contrast, varies only mildly ($F = 0.39$--$0.85$) across the whole taxonomy, consistent with the SFT result that interference is a property of the protocol more than of the teacher. Gaussian universality, in short, is a good guide for some paradigms (SFT) and fails badly for others (RLHF, weak-to-strong).

\section{Discussion and limitations}
\label{sec:discussion}

\paragraph{What kind of ``phase transitions'' are these?}
At $d = 32$--$64$ we observe sharp crossovers, not certified non-analyticities; establishing genuine transitions requires the finite-size scaling limit $d \to \infty$ at fixed $\alpha$, which our framework supports but which we have not pushed to asymptopia for the multi-stage paradigms. We deliberately report boundary \emph{locations} and their trends (e.g.\ $\KL^*(\alpha_r)$ saturating; the transfer sign boundary $\rho^*(\alpha_{\rm ft})$), which are meaningful at finite size and are the quantities a mean-field theory would predict first. The analogy we intend is precise: as in early spin-glass simulations, the measured boundaries constrain future theory regardless of whether they sharpen into true transitions.

\paragraph{Scope of the minimal models.}
Our RLHF policy is a mean-shifted Gaussian, the mildest policy class that can leave the data distribution; richer policy classes could hack the proxy at even smaller KL, so our $\KL^*$ is best read as an upper bound. The SFT similarity dial interpolates weights, which conflates input-feature and readout similarity; the ICL transformer is far from LLM scale; and the trained-teacher family currently contains one dataset (digits). Each limitation is an axis of the released framework rather than a wall: teachers, input distributions, and paradigms are injected through one interface, so e.g.\ a CIFAR-trained or language-model teacher drops into every experiment of this paper unchanged.

\paragraph{Why measure before deriving?}
One may object that numerics without theory cannot explain. Our position is that the statistical-mechanics program has always alternated between the two: solvable models tell us \emph{which} observables are macroscopic, and experiments on harder systems tell theory \emph{where} the interesting structure is. The taxonomy cross is the clearest example---the $40\times$ spread in $\KL^*$ across teacher ensembles is a concrete, falsifiable regularity that no current analysis predicts, and it was cheap to find once the paradigms were instrumented with the right order parameters.

\section{Related work}
\label{sec:related}

\paragraph{Solvable teacher--student models.}
The replica and online analyses of perceptrons and committee machines \citep{gardner1989three,seung1992statistical,saad1995exact,engel2001statistical} established order parameters and phase transitions as the language of learning theory; modern extensions cover structured data \citep{goldt2020modelling,gerace2020generalisation}, kernel and random-feature limits \citep{mei2022generalization,loureiro2021learning}, and SGD dynamics via DMFT \citep{mignacco2020dynamical}; see \citet{zdeborova2016statistical} for a review. We inherit the observables of this line but replace exact solvability with dense numerics on multi-stage nonlinear protocols this theory has not reached.

\paragraph{Empirical phenomenology of modern paradigms.}
Reward overoptimization was quantified at LLM scale by \citet{gao2023scaling}; weak-to-strong generalization by \citet{burns2023weak}; model collapse by \citet{shumailov2024ai}; the task-diversity origin of in-context learning by \citet{garg2022can} and \citet{raventos2023pretraining}; catastrophic forgetting is classical \citep{french1999catastrophic,kirkpatrick2017overcoming}. Relative to these, our contribution is not new phenomena but a \emph{common minimal frame}: all five paradigms reduced to teacher--student form, measured with shared order parameters, on a shared teacher taxonomy---making the phenomena comparable, sweepable, and cheap to reproduce.

\paragraph{Emergence and transitions in deep learning.}
Grokking \citep{power2022grokking}, deep double descent \citep{nakkiran2021deep}, and emergent abilities debates connect abrupt capability changes to underlying continuous order parameters; our framework treats the modern paradigms in the same spirit, and its solvable-baseline studies (committee specialization, double descent) validate the numerics against known theory.

\section{Conclusions}
\label{sec:conclusion}

We introduced a teacher--student atlas of five modern learning paradigms---supervised fine-tuning, reward-model optimization, weak-to-strong supervision, recursive synthetic-data training, and in-context learning---formulated minimally enough for dense parameter sweeps yet nonlinear enough to retain each paradigm's characteristic phenomenon. Measured through physics-style order parameters, every paradigm exhibited sharp empirical phase structure: a transfer--forgetting boundary, a data-dependent Goodhart point, a supervisor-competence window for weak-to-strong gains, a small real-data anchor arresting collapse, and a task-diversity transition to in-context regression. A nine-teacher taxonomy showed these boundaries shift systematically---sometimes by more than an order of magnitude---with teacher structure, delimiting where Gaussian universality can and cannot be trusted.

We see this atlas as the experimental half of a program whose theoretical half is actively advancing: as exact analyses extend toward multi-stage and preference-based learning, the boundaries measured here are the quantities they should predict. The released framework is built for that dialogue---new teachers, data manifolds, and paradigms compose through one interface---and we hope it serves as a shared bench where minimal theory and modern practice meet.

\bibliographystyle{unsrtnat}
\bibliography{references}

\appendix

\section{Complete experimental settings}
\label{app:settings}

\paragraph{Shared toolbox.} Students are MLPs $f(\bx) = \boldsymbol{a}^\top \tanh(W \bx/\sqrt{d})$ with widths in the main text; training is full-batch Adam (lr $5\times10^{-3}$, up to $2000$--$3000$ epochs, early stopping patience $20$ except where trajectories are tracked); label noise $\sigma = 0.1$ on real labels. Probe sets are fresh i.i.d.\ samples from the paradigm's input distribution.

\paragraph{SFT.} $d=64$, widths $16/16$, $\alpha_{\rm pre}=8$; grid $\rho \in \{0, 0.25, 0.5, 0.75, 0.9, 1.0\}$, $\alpha_{\rm ft} \in \{0.5, 1, 2, 4, 8, 16\}$; $20$ checkpoints $\times$ $100$ epochs; $3$ seeds.

\paragraph{RLHF.} $d=32$, width $16$; gold reward $r^*(\bx) = f^*(\bx) - \lambda \max(0, \|\bx\|^2/d - \tau)^2$ with $\lambda = 2, \tau = 1$; Bradley--Terry sharpness $\beta = 5$; $\alpha_r \in \{1,2,4,8,16,32,64\}$; KL-coefficient grid $c \in [0.01, 3]$ (nine points, log-spaced), $1200$ policy steps; $3$ seeds. $\KL(\pi_\bmu \| \pi_0) = \|\bmu\|^2/2$ exactly for the Gaussian base distribution.

\paragraph{Weak-to-strong.} $d=64$; widths: truth $16$, weak $2$, strong $32$; grid $\alpha_w \in \{0.5,1,2,4,8,10\}$, $\alpha_s \in \{1,\dots,32\}$; $5$ seeds for the scatter, $2$ for the PGR surface.

\paragraph{Collapse.} $d=64$, width $16$, $\alpha = 8$, $12$ generations, $p_{\rm real} \in \{0, 0.05, 0.1, 0.2, 0.3, 0.5, 1\}$; $5$ seeds.

\paragraph{ICL.} Causal transformer: $4$ layers, $4$ heads, width $128$; $d=8$, context $k=16$; pretraining $3000$ steps of Adam on sequences from a pool of $N_{\rm tasks}$ regression vectors; evaluation on $2048$ pool and unseen sequences; $2$ seeds.

\paragraph{Taxonomy cross.} Teachers as in Section~\ref{sec:set-tax} ($d = 64$; $32$ for RLHF; the digit teachers use the $1797$-image sklearn digits set, globally standardized, with jitter $0.1$). Probe conditions per paradigm as in Figure~\ref{fig:taxonomy}; budgets: $2000$ training epochs, $2$ seeds per cell.

\section{Reproducibility}
Every figure maps to one command of the released framework (\texttt{statphys study \emph{name}}); each command writes the figure and the raw JSON of every panel. Tests cover all protocols and the taxonomy registry; the full suite runs in $\sim$$6$ minutes on a CPU node.

\end{document}
